\section{Prototypes}

\begin{frame}{Prototypes}

\begin{itemize}
\item Concepts are organized in groups around a \txx{prototype}
\item These have typical members (remembered as \txx{exemplars}) 
  \begin{itemize}
  \item What is typical \con{furniture}?
  \item What is a typical \con{bird}?
  \end{itemize}
\item prototypes have \txx{characteristic feature}s: \\a \con{bird}
  \begin{itemize}
  \item has feathers
  \item warbles
  \item flies
  \item lays eggs
  \end{itemize}
\item This work was pioneered by \citep{Rosch:1973,Rosch:Mervis:Gray:Johnson:Boyes-Braem:1976} (very readable)
\end{itemize}

 \end{frame} 

%\begin{frame}{What is a protoypical bird?}

% \\includegraphics[width=\textwidth]{pics/prot

%\end{frame}

\begin{frame}{Relations between Concepts}

\begin{itemize}
\item Concepts are linked in many ways
\item Most common relationship is \txx{hypernymy}: \con{dog} is-a \con{animal}
\item Typically subordinate terms inherit from superordinate terms
\item Larger units of knowledge, such as \txx{frames} are similar
\item Much recent computational work on these
  \begin{itemize}
  \item WordNet
  \item FrameNet
  \end{itemize}
\end{itemize}



\end{frame}

\begin{frame}{Basic Level Categories}
\begin{itemize}
\item Some categories (concepts) seem to be more psychologically basic than others
  \begin{itemize}
  \item Pictures of objects are categorized faster at the basic level
  \item Basic level names used more often in free-naming tasks
  \item Children learn them earlier
  \item Basic-level names are more common in adult discourse 
  \item Basic-level categories are common in different cultures
  \item Basic level names tend to be short
  \item Basic-level names tend to be common in compound nouns
  \end{itemize}
\item
  \begin{tabular}[t]{lll}
superordinate & basic & subordinate \\  \hline
\lex{vehicle} & \lex{bus} & \lex{school bus} \\
\lex{jewelry} & \lex{necklace} & \lex{bib necklace} \\
\lex{animal} & \lex{dog} & \lex{poodle} 
\end{tabular}
\end{itemize}

\end{frame}

\begin{frame}{}
\begin{itemize}
\item Basic level categories are a decomposition of the world into
  maximally informative categories. 
  \begin{itemize}
  \item BLCs maximize the number of attributes shared by members of
    the category
  \item BLCs minimize the number of attributes shared with
    other categories
  \end{itemize}
\item It can be hard to agree on what is the Basic Level: whereas dog
  as a basic category is a species, bird or fish are at a higher
  level, etc.  The levels may be different for different groups of people.
\item Similarly, the notion of frequency is very closely tied
  to the basic level, but not exactly the same.
\end{itemize}

\end{frame}

\begin{frame}{Linguistic Relativity}

\begin{itemize}
\item The language we think in makes some concepts easy to express,
  and some concepts hard
\item The idea behind \txx{linguistic relativity} is that this will
  effect how you think
  \begin{itemize}
  \item Korean lexicalizes politeness and has rigid social hierarchies
  \item English and Chinese speakers differ as to whether they
    concepualize things as substances or individuals
  \item Gendered language speakers have different connotations: \lex{key}
    \begin{itemize}
    \item (German: masculine) `hard, heavy, jagged, metal, and useful'
    \item (Spanish: feminine)  `golden, intricate, little, lovely, shiny, and tiny'
    \end{itemize}
  \item It is easier to differentiate colors that you have names for
  \end{itemize}
\item Most confirmed differences are very, very subtle
\end{itemize}


\end{frame}

\begin{frame}{The Sapir-Whorf Hypothesis}

\begin{description}
\item[strong]  language determines thought and  linguistic categories limit and determine cognitive categories
\item[weak]  linguistic categories and usage influence thought and certain kinds of non-linguistic behaviour
\end{description}

The terms "Strong/Weak Sapir-Whorf Hypothesis" are widely used even though Edward
Sapir and Benjamin Lee Whorf never co-authored anything and never
stated their ideas in terms of a hypothesis let alone with two versions.

\end{frame}

\begin{frame}{What Whorf actually said}
\begin{quotation}
  We dissect nature along lines laid down by our native language. The
  categories and types that we isolate from the world of phenomena we
  do not find there because they stare every observer in the face; on
  the contrary, the world is presented in a kaleidoscope flux of
  impressions which has to be organized by our minds—and this means
  largely by the linguistic systems of our minds. We cut nature up,
  organize it into concepts, and ascribe significances as we do,
  largely because we are parties to an agreement to organize it in
  this way—an agreement that holds throughout our speech community and
  is codified in the patterns of our language [\ldots] all observers
  are not led by the same physical evidence to the same picture of the
  universe, unless their linguistic backgrounds are similar, or can in
  some way be calibrated.
\\ \hfill Whorf (Carroll; Ed.); 1956: pp. 212–214
\end{quotation}
\end{frame}